\chapter{Sistemas LIT y la convoluci\'{o}n}

\section{Introducci\'{o}n}

Hemos visto que la integral de convoluci\'{o}n 
\begin{equation}
        y\left( t\right) =
        \int_{-\infty }^{\infty }x\left( \tau \right) h\left(t-\tau \right) d\tau =
        x\left( t\right) * h\left( t\right) ,
\end{equation}
define la respuesta $y(t)$ de un sistema $LIT$ para una entrada $x(t)$ si conocemos la respuesta al 
impulso $h(t)$. Por lo tanto, en este cap\'{\i}tulo usaremos la convoluci\'{o}n como herramienta para 
estudiar distintas propiedades de los sistemas LIT.

\section{Propiedades de los sistemas LIT}

\subsection{La propiedad conmutativa}

El resultado de la convoluci\'{o}n es independiente del ordenamiento de las se\~{n}ales, o en otras 
palabras, un sistema LIT de respuesta al impulso $h(t)$ excitado con una entrada $x(t)$ genera una salida 
similar al de un sistema LIT de respuesta al impulso $x(t)$ excitado con una entrada $h(t)$.

Consideremos la operaci\'{o}n de convoluci\'{o}n
\begin{equation}
        y\left( t\right) =
                \int_{-\infty }^{\infty }x\left( \tau \right) h\left(t-\tau \right) d\tau =
                \int_{-\infty }^{\infty }h\left( \tau \right) x\left(t-\tau \right) d\tau .
\end{equation}

Realizamos un cambio de variable $\upsilon =t-\tau $ y por consecuencia definimos $\tau =t-\upsilon $ 
para obtener
\begin{equation}
        y\left( t\right) = 
                -\int_{\infty }^{-\infty }x\left( t-\upsilon \right)h\left( \upsilon \right) d\upsilon =
                \int_{-\infty }^{\infty }x\left(t-\upsilon \right) h\left( \upsilon \right) d\upsilon .
\end{equation}

\subsection{La propiedad distributiva}

La salida de un sistema LIT cuya respuesta al impulso es $h_{a} + h_{b}$ para una entrada $x$ es 
similar a la suma de las salidas de dos sistemas de respuestas al impulso $h_{a}$ y $h_{b}$ frente a una 
entrada $x$. Se puede probar la distributividad de la convoluci\'{o}n respecto a la suma de se\~{n}ales 
\begin{eqnarray}
y\left( t\right) &=&x\left( t\right) *\left( h_a\left( t\right)
+h_{b}\left( t\right) \right) , \\
y\left( t\right) &=&x\left( t\right) *h_a\left( t\right) +x\left( t\right)
*h_{b}\left( t\right),
\end{eqnarray}
pues para la integral vale la siguiente propiedad
\begin{equation}
\int_{a}^{b} x(v) \, (f(v) + g(v)) \, dv = 
\int_{a}^{b} x(v) \, f(v) \, dv + \int_{a}^{b} x(v) \, g(v)) \, dv.
\end{equation}

\subsection{La propiedad asociativa}

Queda como ejercicio para el alumno la prueba de esta propiedad.

\begin{eqnarray}
y\left( t\right) &=&x\left( t\right) *\left( h_a\left( t\right)
*h_{b}\left( t\right) \right) , \\
y\left( t\right) &=&\left( x\left( t\right) *h_a\left( t\right) \right)
*h_{b}\left( t\right) , \\
y\left( t\right) &=&x\left( t\right) *h_a\left( t\right) *h_{b}\left(
t\right) .
\end{eqnarray}

\subsection{Preguntas}

\begin{enumerate}
\item \textquestiondown Qu\'{e} significado tiene la propiedad conmutativa de los sistemas lineales?
\item \textquestiondown Qu\'{e} significado tiene la propiedad distributiva de los sistemas lineales?
\item \textquestiondown Qu\'{e} significado tiene la propiedad asociativa de los sistemas lineales?
\item D\'{e} alg\'{u}n ejemplo concreto de estas tres propiedades.
\end{enumerate}

\subsection{Sistemas LIT con y sin memoria}

Un sistema no tiene memoria si su salida en cualquier momento depende
solamente de la entrada de ese momento.

Si un sistema no tiene memoria entonces debe ser $y\left( t\right) =0$, para todo $t\neq 0$, y por o tanto
\begin{equation}
h\left( t\right) = K \, \delta \left( t\right) .
\end{equation}
Si un sistema tiene una respuesta al impulso $h\left( t\right) $ que no es
id\'{e}nticamente cero para $t\neq 0,$ entonces el sistema tiene memoria.

\subsection{Preguntas}

\begin{enumerate}
\item \textquestiondown Conoce alg\'{u}n sistema f\'{\i}sico que no tenga memoria?
\item \textquestiondown Conoce alg\'{u}n sistema f\'{\i}sico que tenga memoria?
\end{enumerate}

\subsection{Sistemas LIT causales y no causales}

Si $y(t)$ no depende de $x(t_{1})$ para $t_{1}>t,$ entonces un sistema LIT es causal y la funci\'{o}n 
$h\left( t-\tau \right) $ que multiplica a $x(t)$ debe ser cero para $t<\tau $, o sea que 
\begin{equation}
h\left( t\right) =0, \, t<0,
\end{equation}
y la integral de convoluci\'{o}n para un sistema LIT causal es 
\begin{equation}
y\left( t\right) =\int_{0}^{t}x\left( \tau \right) \, h\left( t-\tau \right)
d\tau .
\end{equation}

\subsection{Sistemas LIT estables e inestables}

\subsubsection{Estabilidad BIBO}

El criterio de estabilidad \textbf{B}ounded \textbf{I}nput \textbf{B}ounded \textbf{O}utput 
(entrada acotada, salida acotada) se puede expresar matem\'{a}ticamente como 
\begin{equation}
        \left| x\left( t\right) \right| < M_{x} \rightarrow \left| y\left( t\right)\right| <M_{y},
\end{equation}
para tiempo continuo.

El sistema es estable si para entradas acotadas 
\begin{equation}
        \left| \int_{-\infty}^{t}x\left( \tau \right) d\tau \right| < M_{x}.
\end{equation}
tenemos salidas acotadas 
\begin{equation}
        \left| y\left( t\right) \right| =
            \left| \int_{-\infty }^{t}x\left( \tau\right) h\left( t-\tau \right) d\tau \right| <
            \left| \int_{-\infty}^{t}x\left( \tau \right) d\tau \right| 
            \left| \int_{-\infty }^{t}h\left(t-\tau \right) d\tau \right| .
\end{equation}
Reemplazando $\left| \int_{-\infty}^{t}x\left( \tau \right) d\tau \right|$ por la cota $M_{x}$
\begin{equation}
    \left| y\left( t\right) \right| < M_{x}\left| \int_{-\infty }^{t}h\left(t-\tau \right) d\tau \right| = M_{y}
\end{equation}
obtenemos que debe existir una cota
\begin{equation}
       \left| \int_{-\infty }^{t}h\left( t-\tau \right) d\tau \right| < M_{h},
\end{equation}
como condición para la estabilidad BIBO del sistema LIT.

\subsection{Preguntas}

\begin{enumerate}
\item \textquestiondown Conoce alg\'{u}n sistema f\'{\i}sico inestable?
\item De ejemplos de sistemas estables e inestables.
\item Indicar si las siguientes afirmaciones son verdaderas o falsas, en el caso de un sistema lineal 
invariante en el tiempo:
\begin{enumerate}
\item  Si la entrada es nula para todo tiempo, entonces la salida es nula para todo tiempo.

\item  Si la relaci\'{o}n entre entrada y salida es 
$y\left(t\right) =x\left( t\right) \cos \left( t\right) ,$ entonces el sistema  es LTI. 

\item  Si la relaci\'{o}n entre entrada y salida es 
$y\left(t\right) =\int_{-\infty }^{t}x\left( \tau \right) d\tau ,$ entonces el sistema es LTI. 
\end{enumerate}
\end{enumerate}

% \subsection{Respuestas}

% Estas son las respuestas a las preguntas planteadas arriba

% \begin{enumerate}
%     \item  \textbf{Verdadero, pues si la entrada es nula, la integral (o suma) de convoluci\'{o}n es nula.}

%     \item   \textbf{Falso, porque la respuesta al impulso depende del tiempo en que el impulso ocurre.}
%         \begin{equation}
%                 y(t) = \delta(t_{0}) \cos \left( t_{0} \right).
%         \end{equation}

%     \item  \textbf{Verdadero, pues dicha relaci\'{o}n puede escribirse como una convoluci\'{o}n} 
%         \begin{equation}
%                 y\left( t\right) =\int_{-\infty }^{t}x\left( \tau \right) d\tau
%                 =\int_{-\infty }^{t}x\left( \tau \right) u\left( t-\tau \right) d\tau
%                 =x\left( t\right) * u\left( t\right).
%         \end{equation}
% \end{enumerate}
